
\subsection{Corrections des Exercices}

\begin{correctionbox}[Correction Exercice 1 : Chebyshev (Calcul de base)]
$P(|Y - \mu| \ge k) \le \frac{\sigma^2}{k^2}$.
$P(|Y - 50| \ge 10) \le \frac{16}{10^2} = \frac{16}{100} = 0.16$.
La probabilité est d'au maximum 16\%.
\end{correctionbox}

\begin{correctionbox}[Correction Exercice 2 : Chebyshev (En termes d'écarts-types)]
On pose $k = 5\sigma$.
$P(|Y - \mu| \ge 5\sigma) \le \frac{\sigma^2}{(5\sigma)^2} = \frac{\sigma^2}{25\sigma^2} = \frac{1}{25} = 0.04$.
La probabilité est d'au maximum 4\%.
\end{correctionbox}

\begin{correctionbox}[Correction Exercice 3 : Chebyshev (Application à $\bar{X}_n$)]
1.  $E[\bar{X}_{100}] = \mu = 100$.
    $\text{Var}(\bar{X}_{100}) = \frac{\sigma^2}{n} = \frac{400}{100} = 4$.
2.  On applique Chebyshev à $\bar{X}_{100}$ (avec $\mu=100, \sigma^2=4$) et $k=4$.
    $P(|\bar{X}_{100} - 100| \ge 4) \le \frac{\text{Var}(\bar{X}_{100})}{k^2} = \frac{4}{4^2} = \frac{4}{16} = 0.25$.
\end{correctionbox}

\begin{correctionbox}[Correction Exercice 4 : CLT (Distribution de $\bar{X}_n$)]
$E[\bar{X}_{64}] = \mu = 20$.
$\text{Var}(\bar{X}_{64}) = \frac{\sigma^2}{n} = \frac{16}{64} = 0.25$.
Selon le TCL, $\bar{X}_{64} \approx \mathcal{N}(20, 0.25)$.
\end{correctionbox}

\begin{correctionbox}[Correction Exercice 5 : CLT (Calcul de probabilité pour $\bar{X}_n$)]
On standardise $\bar{X}_{64} \approx \mathcal{N}(20, 0.25)$, donc $\sigma_{\bar{X}} = \sqrt{0.25} = 0.5$.
On cherche $P(\bar{X}_{64} \le 21)$.
$Z = \frac{\bar{X}_n - \mu}{\sigma_{\bar{X}}} = \frac{21 - 20}{0.5} = \frac{1}{0.5} = 2$.
$P(\bar{X}_{64} \le 21) = P(Z \le 2) = \Phi(2) \approx 0.9772$.
\end{correctionbox}

\begin{correctionbox}[Correction Exercice 6 : CLT (Calcul de probabilité pour $\bar{X}_n$)]
$\mu=50, \sigma=5, n=100$. $\bar{X}_{100} \approx \mathcal{N}(\mu, \sigma^2/n)$.
$\mu_{\bar{X}} = 50$. $\sigma_{\bar{X}} = \sigma / \sqrt{n} = 5 / \sqrt{100} = 0.5$.
On cherche $P(49 \le \bar{X}_{100} \le 51)$.
$Z_1 = \frac{49 - 50}{0.5} = -2$. $Z_2 = \frac{51 - 50}{0.5} = 2$.
$P(-2 \le Z \le 2) = \Phi(2) - \Phi(-2) = \Phi(2) - (1 - \Phi(2)) = 2\Phi(2) - 1$.
$P \approx 2(0.9772) - 1 = 1.9544 - 1 = 0.9544$. (La règle des 95\%).
\end{correctionbox}

\begin{correctionbox}[Correction Exercice 7 : CLT (Inverse pour $\bar{X}_n$)]
$\mu=10, \sigma^2=81, n=36$. $\bar{X}_{36} \approx \mathcal{N}(\mu, \sigma^2/n)$.
$\mu_{\bar{X}} = 10$. $\sigma_{\bar{X}} = \sigma / \sqrt{n} = 9 / \sqrt{36} = 9 / 6 = 1.5$.
On cherche $c$ tel que $P(\bar{X}_{36} \le c) \approx 0.9772$.
$P(Z \le \frac{c - 10}{1.5}) = 0.9772$.
D'après la table, le $z$ correspondant est $2.0$.
$\frac{c - 10}{1.5} = 2 \implies c - 10 = 3 \implies c = 13$.
\end{correctionbox}

\begin{correctionbox}[Correction Exercice 8 : CLT (Distribution de $S_n$)]
$E[S_{100}] = n\mu = 100 \times 10 = 1000$.
$\text{Var}(S_{100}) = n\sigma^2 = 100 \times 4 = 400$.
Selon le TCL, $S_{100} \approx \mathcal{N}(1000, 400)$.
\end{correctionbox}

\begin{correctionbox}[Correction Exercice 9 : CLT (Calcul de probabilité pour $S_n$)]
On standardise $S_{100} \approx \mathcal{N}(1000, 400)$, donc $\sigma_{S_n} = \sqrt{400} = 20$.
On cherche $P(S_{100} > 1020)$.
$Z = \frac{S_n - n\mu}{\sigma_{S_n}} = \frac{1020 - 1000}{20} = \frac{20}{20} = 1$.
$P(S_{100} > 1020) = P(Z > 1) = 1 - \Phi(1) \approx 1 - 0.8413 = 0.1587$.
\end{correctionbox}

\begin{correctionbox}[Correction Exercice 10 : CLT (Application $S_n$)]
$n=49, \mu=70, \sigma=14$. $S_{49} \approx \mathcal{N}(n\mu, n\sigma^2)$.
$E[S_{49}] = n\mu = 49 \times 70 = 3430$.
$\sigma_{S_n} = \sigma \sqrt{n} = 14 \times \sqrt{49} = 14 \times 7 = 98$.
On cherche $P(S_{49} > 3500)$.
$Z = \frac{3500 - 3430}{98} = \frac{70}{98} \approx 0.714$.
$P(Z > 0.714)$. (Valeur non fournie, mais le calcul est le Z-score).
\end{correctionbox}

\begin{correctionbox}[Correction Exercice 11 : CLT (Inverse pour $S_n$)]
$n=64, \mu=5, \sigma^2=1$. $S_{64} \approx \mathcal{N}(n\mu, n\sigma^2)$.
$E[S_{64}] = 64 \times 5 = 320$.
$\sigma_{S_n} = \sigma \sqrt{n} = 1 \times \sqrt{64} = 8$.
On cherche $c$ tel que $P(S_{64} > c) \approx 0.05$.
$P(Z > \frac{c - 320}{8}) = 0.05$.
$P(Z \le z) = 0.95$. D'après la table, $z \approx 1.645$.
$\frac{c - 320}{8} = 1.645 \implies c = 320 + 8(1.645) = 320 + 13.16 = 333.16$.
\end{correctionbox}

\begin{correctionbox}[Correction Exercice 12 : Règle d'Approximation]
$n=40, p=0.1$.
$np = 40 \times 0.1 = 4$.
$n(1-p) = 40 \times 0.9 = 36$.
Puisque $np=4$ est $< 10$, la règle n'est pas satisfaite. L'approximation normale n'est pas recommandée.
\end{correctionbox}

\begin{correctionbox}[Correction Exercice 13 : Règle d'Approximation]
$n=500, p=0.04$.
$np = 500 \times 0.04 = 20$.
$n(1-p) = 500 \times 0.96 = 480$.
Les deux conditions ($20 \ge 10$ et $480 \ge 10$) sont satisfaites. L'approximation est valide.
\end{correctionbox}

\begin{correctionbox}[Correction Exercice 14 : Paramètres d'Approximation]
$X \sim \text{Bin}(100, 0.3)$.
$\mu_Y = np = 100 \times 0.3 = 30$.
$\sigma_Y^2 = np(1-p) = 100 \times 0.3 \times 0.7 = 21$.
L'approximation est $Y \sim \mathcal{N}(30, 21)$.
\end{correctionbox}

\begin{correctionbox}[Correction Exercice 15 : Correction de Continuité (Règles)]
1. $P(X = 50) \approx P(49.5 \le Y \le 50.5)$
2. $P(X \ge 50) \approx P(Y \ge 49.5)$
3. $P(X \le 49) \approx P(Y \le 49.5)$
4. $P(X < 49) = P(X \le 48) \approx P(Y \le 48.5)$
\end{correctionbox}

\begin{correctionbox}[Correction Exercice 16 : Correction de Continuité (Règles)]
1. $P(X > 30) = P(X \ge 31) \approx P(Y \ge 30.5)$
2. $P(30 < X < 40) = P(31 \le X \le 39) \approx P(30.5 \le Y \le 39.5)$
3. $P(30 \le X \le 40) \approx P(29.5 \le Y \le 40.5)$
\end{correctionbox}

\begin{correctionbox}[Correction Exercice 17 : Binomiale (Calcul P(X=k))]
$X \sim \text{Bin}(100, 0.5)$. $\mu = 50$, $\sigma^2 = 25$, $\sigma = 5$.
On cherche $P(X=50) \approx P(49.5 \le Y \le 50.5)$.
$Z_1 = \frac{49.5 - 50}{5} = -0.1$. $Z_2 = \frac{50.5 - 50}{5} = 0.1$.
$P(-0.1 \le Z \le 0.1) = \Phi(0.1) - \Phi(-0.1) = \Phi(0.1) - (1 - \Phi(0.1)) = 2\Phi(0.1) - 1$.
$P \approx 2(0.5398) - 1 = 1.0796 - 1 = 0.0796$.
\end{correctionbox}

\begin{correctionbox}[Correction Exercice 18 : Binomiale (Calcul $P(X \ge k)$)]
$X \sim \text{Bin}(100, 0.5)$. $\mu = 50$, $\sigma = 5$.
On cherche $P(X \ge 60) \approx P(Y \ge 59.5)$.
$Z = \frac{59.5 - 50}{5} = \frac{9.5}{5} = 1.9$.
$P(Z \ge 1.9) = 1 - \Phi(1.9)$. (Valeur $\Phi(1.9)$ non fournie, mais $1-\Phi(1.96) \approx 0.025$).
\end{correctionbox}

\begin{correctionbox}[Correction Exercice 19 : Binomiale (Calcul $P(X \le k)$)]
$X \sim \text{Bin}(400, 0.2)$. $np = 80$, $n(1-p)=320$. (Règle OK).
$\mu = 80$. $\sigma^2 = 80 \times 0.8 = 64$. $\sigma = 8$.
On cherche $P(X \le 70) \approx P(Y \le 70.5)$.
$Z = \frac{70.5 - 80}{8} = \frac{-9.5}{8} \approx -1.19$.
$P(Z \le -1.19) = 1 - \Phi(1.19)$. (Valeur non fournie).
\end{correctionbox}

\begin{correctionbox}[Correction Exercice 20 : Binomiale (Calcul $P(X > k)$)]
$X \sim \text{Bin}(400, 0.2)$. $\mu = 80$, $\sigma = 8$.
On cherche $P(X > 92) = P(X \ge 93) \approx P(Y \ge 92.5)$.
$Z = \frac{92.5 - 80}{8} = \frac{12.5}{8} = 1.5625$.
$P(Z \ge 1.5625) \approx P(Z \ge 1.58) = 1 - \Phi(1.58) \approx 1 - 0.9429 = 0.0571$.
\end{correctionbox}

\begin{correctionbox}[Correction Exercice 21 : Binomiale (Calcul $P(k_1 \le X \le k_2)$)]
$X \sim \text{Bin}(100, 0.6)$. $np = 60$, $n(1-p)=40$. (Règle OK).
$\mu = 60$. $\sigma^2 = 60 \times 0.4 = 24$. $\sigma = \sqrt{24} \approx 4.9$.
On cherche $P(50 \le X \le 65) \approx P(49.5 \le Y \le 65.5)$.
$Z_1 = \frac{49.5 - 60}{4.9} = \frac{-10.5}{4.9} \approx -2.14$.
$Z_2 = \frac{65.5 - 60}{4.9} = \frac{5.5}{4.9} \approx 1.12$.
$P(-2.14 \le Z \le 1.12) = \Phi(1.12) - \Phi(-2.14) = \Phi(1.12) - (1 - \Phi(2.14))$. (Valeurs non fournies).
\end{correctionbox}

\begin{correctionbox}[Correction Exercice 22 : TCL vs Chebyshev]
$n=100, \mu=10, \sigma^2=25$. $\bar{X}_{100} \approx \mathcal{N}(10, 25/100=0.25)$. $\sigma_{\bar{X}} = 0.5$.
On cherche $P(|\bar{X}_{100} - 10| \ge 1)$.
1.  **Chebyshev** : $k=1, \text{Var}(\bar{X}_{100}) = 0.25$.
    $P \le \frac{\text{Var}(\bar{X}_{100})}{k^2} = \frac{0.25}{1^2} = 0.25$. (Borne $\le 25\%$).
2.  **CLT** : $P(|\bar{X}_{100} - 10| \ge 1) = P(Z \ge \frac{1}{0.5}) + P(Z \le \frac{-1}{0.5}) = P(Z \ge 2) + P(Z \le -2)$.
    $P = (1 - \Phi(2)) + \Phi(-2) = (1 - \Phi(2)) + (1 - \Phi(2)) = 2(1 - \Phi(2))$.
    $P \approx 2(1 - 0.9772) = 2(0.0228) = 0.0456$. (Probabilité $\approx 4.56\%$).
\end{correctionbox}

\begin{correctionbox}[Correction Exercice 23 : Taille d'Échantillon (CLT)]
$\bar{X}_n \approx \mathcal{N}(\mu, \sigma^2/n) = \mathcal{N}(0, 1/n)$. $\sigma_{\bar{X}} = 1/\sqrt{n}$.
On veut $P(|\bar{X}_n| \le 0.1) \ge 0.95$, soit $P(-0.1 \le \bar{X}_n \le 0.1) \ge 0.95$.
On standardise : $P\left( \frac{-0.1 - 0}{1/\sqrt{n}} \le Z \le \frac{0.1 - 0}{1/\sqrt{n}} \right) \ge 0.95$.
$P(-0.1\sqrt{n} \le Z \le 0.1\sqrt{n}) \ge 0.95$.
On sait que $P(-1.96 \le Z \le 1.96) = 0.95$.
On doit donc avoir $0.1\sqrt{n} \ge 1.96$.
$\sqrt{n} \ge 19.6 \implies n \ge (19.6)^2 = 384.16$.
Il faut $n=385$ échantillons au minimum.
\end{correctionbox}

\begin{correctionbox}[Correction Exercice 24 : LLN vs CLT]
$\mu=10, \sigma^2=100$.
1.  **$n=100$** : $\bar{X}_{100} \approx \mathcal{N}(10, 100/100=1)$. $\sigma_{\bar{X}} = 1$.
    $P(9 \le \bar{X}_{100} \le 11) = P(\frac{9-10}{1} \le Z \le \frac{11-10}{1}) = P(-1 \le Z \le 1)$.
    $P = 2\Phi(1) - 1 \approx 2(0.8413) - 1 = 0.6826$.
2.  **$n=10000$** : $\bar{X}_{10000} \approx \mathcal{N}(10, 100/10000=0.01)$. $\sigma_{\bar{X}} = 0.1$.
    $P(9 \le \bar{X}_{10000} \le 11) = P(\frac{9-10}{0.1} \le Z \le \frac{11-10}{0.1}) = P(-10 \le Z \le 10)$.
    $P \approx \Phi(10) - \Phi(-10) \approx 1 - (1-1) = 1$.
3.  **Illustration** : Le CLT montre \textit{comment} la LLN fonctionne. En augmentant $n$, la variance de $\bar{X}_n$ s'effondre (de 1 à 0.01), forçant la distribution de $\bar{X}_n$ à se concentrer massivement autour de $\mu=10$. La probabilité que $\bar{X}_n$ soit proche de $\mu$ tend vers 1.
\end{correctionbox}

\begin{correctionbox}[Correction Exercice 25 : Binomiale (Sans Correction)]
$X \sim \text{Bin}(400, 0.1)$. On approxime $P(X \le 30.5)$.
$\mu = np = 40$. $\sigma^2 = np(1-p) = 36$. $\sigma = 6$.
$Y \sim \mathcal{N}(40, 36)$.
On cherche $P(Y \le 30.5)$.
$Z = \frac{30.5 - 40}{6} = \frac{-9.5}{6} \approx -1.583$.
$P(Z \le -1.58) = \Phi(-1.58) = 1 - \Phi(1.58) \approx 1 - 0.9429 = 0.0571$.
(Note : C'est le même calcul que l'exercice 19, $P(X \le 30)$, car $P(X \le 30) \approx P(Y \le 30.5)$).
\end{correctionbox}
